\section*{الفصل \ref{ch:g11:mutations} --- الطفرات والتنوع الوراثي}

\begin{solution}{exo:g11:mutations:1}
تغير في متتالية نوكليوتيدات \omterm{def:g10:universal-dna:information}{DNA}، ينتقل إلى ذرية \omterm{def:g10:cells-common-unit:cell}{الخلية}. والأصناف:
استبدال، وإقحام، وحذف.
\end{solution}

\begin{solution}{exo:g11:mutations:2}
\texttt{ATGCCAGAAT}: استبدال (T بالقاعدة A في الموضع 6).
و\texttt{ATGCCTGAT}: حذف حرف A واحد. و\texttt{ATGCCTTGAAT}: إقحام
حرف T واحد.
\end{solution}

\begin{solution}{exo:g11:mutations:3}
الجسدية: في \omterm{def:g10:cells-common-unit:cell}{خلية} جسدية، لا تنتقل إلا إلى ذرية تلك \omterm{def:g10:cells-common-unit:cell}{الخلية} داخل الفرد.
والسلالية: في مشيج أو سليفه، تنتقل إلى النسل وإلى كل خلاياه. ولا تدخل
في الوراثة إلا \omterm{def:g11:mutations:mutation}{الطفرات} السلالية.
\end{solution}

\begin{solution}{exo:g11:mutations:4}
الضوء فوق البنفسجي، والأشعة السينية، وبنزوبيرين دخان التبغ (وكذلك
النشاط الإشعاعي والأفلاتوكسين). وفوق البنفسجي يلحم تيمينين متجاورين من
شريط واحد، فيشوه اللولب.
\end{solution}

\begin{solution}{exo:g11:mutations:5}
الشريط المتمم السليم: فلما كان كل شريط يحدد الآخر، أمكن إعادة بناء
المتتالية الصحيحة للقطعة المتضررة من شريكها.
\end{solution}

\begin{solution}{exo:g11:mutations:6}
العادية: نحو 6\%؛ والناقصة الإصلاح: 0.01\%. فالإصلاح يحسّن البقاء نحو
600 ضعف عند تلك الجرعة.
\end{solution}

\begin{solution}{exo:g11:mutations:7}
$20/10^9 = \num{2e-8}$. ولا: فالطبق لا يكشف إلا أي \omterm{def:g10:cells-common-unit:cell}{الخلايا} كانت
مقاومة من قبل؛ أما كون المضاد الحيوي سبب \omterm{def:g11:mutations:mutation}{الطفرة} فيحتاج إلى اختبار
النسخ.
\end{solution}

\begin{solution}{exo:g11:mutations:8}
تتلقى كل نسخة \omterm{def:g10:cells-common-unit:cell}{خلايا} من المستعمرات الأصلية نفسها. فلو أن المضاد سبّب
المقاومة عند التماس، لكانت \omterm{def:g10:cells-common-unit:cell}{الخلايا} المصابة قلة عشوائية على كل نسخة، في
مواضع لا صلة بينها. أما وجود المقاومة دائمًا في المواضع نفسها فيعني أن
تلك المستعمرات الأصلية، وقد نمت بلا دواء، كانت تحوي \omterm{def:g10:cells-common-unit:cell}{خلايا} مقاومة من
قبل.
\end{solution}

\begin{solution}{exo:g11:mutations:9}
الأخطاء لكل نسخة من \omterm{def:g10:universal-dna:gene}{المورّثة}: $2 \times 3000 \times 10^{-9} =
\num{6e-6}$ (على الشريطين). ومن بين $10^{12}$ \omterm{def:g10:cells-common-unit:cell}{خلية} أنتجت في الانقسام
الأخير، نحو $\num{6e6}$ تحمل \omterm{def:g11:mutations:mutation}{طفرة} جديدة في تلك \omterm{def:g10:universal-dna:gene}{المورّثة}.
\end{solution}

\begin{solution}{exo:g11:mutations:10}
\omterm{def:g11:mutations:mutation}{طفرة} جسدية في \omterm{def:g10:cells-common-unit:cell}{خلية} برعم أعطت الغصن لونه؛ والعقل هي ذلك النسيج، فتحفظه.
أما \omterm{def:g10:cells-common-unit:cell}{الخلايا} الجنسية في الأزهار فتشتق من الغصن نفسه، لكن تغيرًا في \omterm{def:g10:universal-dna:gene}{أليل}
واحد لا ينتقل إلا إلى نصف الأمشاج، وقد لقّحت البييضات بحبوب طلع تحمل
\omterm{def:g10:universal-dna:gene}{الأليل} العادي؛ \omterm{def:g10:universal-dna:gene}{فالأليل} الجديد، إن كان متنحيًا، يختفي في البادرات.
\end{solution}

\begin{solution}{exo:g11:mutations:11}
الإنزيمة الناقصة تصلح ضرر فوق البنفسجي؛ وفوق البنفسجي لا ينفذ إلا إلى
الجلد. أما المعي فلا يتلقى فوق بنفسجي، ولا يصاب بمثل تلك الآفات، ولا
يحتاج إلى ذلك الإصلاح.
\end{solution}

\begin{solution}{exo:g11:mutations:12}
الآفات تتناسب مع الجرعة: فعشر فوق البنفسجي يعني عشر الآفات في \omterm{def:g10:cells-common-unit:cell}{الخلية}
في الساعة --- وهي مع ذلك آلاف. وحروق الشمس موت \omterm{def:g10:cells-common-unit:cell}{خلايا} كثيرة؛ ودون تلك
الجرعة تنجو \omterm{def:g10:cells-common-unit:cell}{الخلايا}، لكن الآفات التي تفلت من الإصلاح تصير \omterm{def:g11:mutations:mutation}{طفرات}. فغياب
الحرق يعني أن الضرر كان محتملًا، لا أنه غائب.
\end{solution}

\begin{solution}{exo:g11:mutations:13}
ما من واحد منهما في ذاته. ففي منطقة فيها ملاريا ترفع النسخة الواحدة
البقاء، ويظل \omterm{def:g10:universal-dna:gene}{الأليل} شائعًا رغم مرض حاملي النسختين؛ وحيث تغيب الملاريا
لا يسبب إلا المرض ويبقى نادرًا. فالمحيط، وعدد النسخ، يقرران قيمة
\omterm{def:g11:mutations:mutation}{الطفرة}.
\end{solution}

\begin{solution}{exo:g11:mutations:14}
تموت \omterm{def:g10:cells-common-unit:cell}{الخلايا} الحساسة؛ وتنقسم المقاومة، وهي موجودة من قبل، بحرية،
فتصير المزرعة كلها مقاومة في غضون يوم. والجرعة التي توقف باكرًا تترك
آخر الناجيات --- وهي الأشد مقاومة --- حية لتنمو من جديد، في جسم صار
غنيًا بالمقاومة؛ أما أخذ الجرعة كاملة فيقتلها قبل أن تتكاثر.
\end{solution}

\begin{solution}{exo:g11:mutations:15}
معظم \omterm{def:g11:mutations:mutation}{الطفرات} محايدة (خارج \omterm{def:g10:universal-dna:gene}{المورّثات} أو بلا أثر في \omterm{def:g10:chemistry-of-life:families}{البروتين})، وأقلية
منها ضارة، وقليل جدًا نافع؛ وهي، لكونها عشوائية، تقع بصرف النظر عن
الفائدة، لكن الانتقاء يبقي النافعة ويزيل الضارة. والفائدة تتوقف على
المحيط، \omterm{def:g11:mutations:mutation}{فالطفرة} الضارة اليوم قد تكون هي المهمة حين تتغير الظروف. وعبر
آلاف الأجيال، ومع مليارات الأفراد، تنشأ حتى \omterm{def:g11:mutations:mutation}{الطفرات} النافعة النادرة
مرارًا: فالتطور مبني على الباقي المفروز، لا على متوسط \omterm{def:g11:mutations:mutation}{الطفرات}.
\end{solution}

\begin{solution}{pb:g11:mutations:1}
\textbf{1.} حتى تكون كل مستعمرة نسيلة معروفة الأصل، وحتى تكون
المستعمرات متباعدة بما يكفي للتمييز بينها على النسخ.

\textbf{2.} بضع \omterm{def:g10:cells-common-unit:cell}{خلايا} من أعلى كل مستعمرة، بالترتيب الهندسي نفسه الذي
على الطبق، لأن الوسادة تضغط مسطحة من دون انزلاق.

\textbf{3.} أن \omterm{def:g10:cells-common-unit:cell}{الخلايا} المقاومة كانت موجودة في تينك المستعمرتين
الأصليتين: فثلاث نسخ مستقلة لا تتفق بالمصادفة.

\textbf{4.} مستعمرات مقاومة في مواضع عشوائية لا صلة بينها على الأطباق
الثلاثة، لأن أي \omterm{def:g10:cells-common-unit:cell}{الخلايا} استُحثت يختلف من نسخة إلى أخرى.

\textbf{5.} \omterm{def:g11:mutations:mutation}{طفرة} قبل التماس: فالمقاومة ترسم على المستعمرات الأصلية،
وهي لم تلق الدواء قط.

\textbf{6.} \omterm{def:g10:cells-common-unit:cell}{الخلايا} من الموضعين المقاومين تنمو على الستربتوميسين؛
\omterm{def:g10:cells-common-unit:cell}{والخلايا} من المواضع الأخرى لا تنمو.

\textbf{7.} بعضها فقط على الأرجح: \omterm{def:g11:mutations:mutation}{فالطفرة} وقعت في انقسام ما أثناء نمو
المستعمرة، ولا يقاوم إلا ذرية تلك \omterm{def:g10:cells-common-unit:cell}{الخلية}. وللبت في ذلك، ينثر عدد معلوم
من \omterm{def:g10:cells-common-unit:cell}{خلايا} المستعمرة على الستربتوميسين وتعد الناجيات.

\textbf{8.} في الانقسام العشرين من ثلاثة وعشرين: $1/2^{20-1}$\dots
وأبسط من ذلك، تنقسم \omterm{def:g10:cells-common-unit:cell}{الخلية} الطافرة 3 مرات أخرى: $2^3 = 8$ \omterm{def:g10:cells-common-unit:cell}{خلايا} مقاومة
من أصل $2^{23}$، أي نحو واحدة في المليون. وفي الانقسام الثالث: واحدة
من \omterm{def:g10:cells-common-unit:cell}{الخلايا} الثماني الموجودة عندئذ، أي ثمن المستعمرة.

\textbf{9.} حتى بضع \omterm{def:g10:cells-common-unit:cell}{خلايا} مقاومة بين نحو $10^4$ \omterm{def:g10:cells-common-unit:cell}{خلية} يرفعها المخمل
تكفي لتأسيس مستعمرة على طبق المضاد؛ فالاختبار يكشف وجود \omterm{def:g10:cells-common-unit:cell}{خلايا} مقاومة،
لا نسبتها.

\textbf{10.} تقوم الحجة على الموضع: فلا يمكن أن تطابَق مستعمرة على
نسخة بمستعمرة على الطبق الأصلي إلا إذا حفظت الهندسة.

\textbf{11.} $30/\num{10000} = \num{3e-3}$.

\textbf{12.} نحو $\num{3e-3}$ \omterm{def:g11:mutations:mutation}{طفرة} لكل $10^7$ انقسام:
أي $\num{3e-10}$ لكل انقسام.

\textbf{13.} من رتبة نسبة الخطأ نفسها عند \omterm{def:g10:universal-dna:nucleotide}{نوكليوتيد} واحد: \omterm{def:g11:mutations:mutation}{فطفرة}
المقاومة في جوهرها استبدال واحد بعينه أفلت من الإصلاح.

\textbf{14.} $10^9$ انقسام عند $\num{3e-10}$: أي وسطيًا كسر من \omterm{def:g11:mutations:mutation}{طفرة}
واحدة، فمعظم المزارع لا تحمل شيئًا أو تحمل بضع \omterm{def:g10:cells-common-unit:cell}{خلايا} مقاومة (تضاعف كل
\omterm{def:g11:mutations:mutation}{طفرة} بعد ذلك الانقسامات التي تليها).

\textbf{15.} \omterm{def:g11:mutations:mutation}{الطفرة} في انقسام باكر يرثها جزء كبير من المزرعة، فيكون
العدد كبيرًا؛ أما المتأخرة فلا تعطي إلا بضع \omterm{def:g10:cells-common-unit:cell}{خلايا} مقاومة. ولأن التوقيت
عشوائي، تعطي المزارع المكررة أعدادًا شديدة التشتت --- وهي نفسها علامة
على أن \omterm{def:g11:mutations:mutation}{الطفرات} تنشأ قبل المضاد الحيوي لا استجابة له.

\textbf{16.} نادرة: أي أن \omterm{def:g11:mutations:mutation}{الطفرات} تقع قليلًا في كل انقسام. وعشوائية:
أي أنها تقع باستقلال عن فائدتها. والتجربة ترسي العشوائية (بحجة
المواضع)؛ أما أعداد الجزء الثالث فتقيس الندرة.

\textbf{17.} فرز الجمهرة: فقد قتل \omterm{def:g10:cells-common-unit:cell}{الخلايا} الحساسة وترك المقاومة
تتكاثر --- أي انتقاءً.

\textbf{18.} وفّرت \omterm{def:g11:mutations:mutation}{الطفرة} \omterm{def:g10:universal-dna:gene}{الأليل} المقاوم، عشوائيًا، قبل أي حاجة؛ وفرز
المضاد الحيوي الجمهرة فلم يبقِ إلا حامليه.

\textbf{19.} \omterm{def:g11:mutations:mutation}{الطفرة} تغير جزيئًا يدخل في فعل الدواء، لا نمو المستعمرة
ولا مظهرها من دون الدواء؛ والفرق الوراثي لا يكون له أثر مرئي إلا في
محيط يهم فيه.

\textbf{20.} أثبتت المواضع المتطابقة على النسخ الثلاث أن \omterm{def:g11:mutations:mutation}{طفرات}
المقاومة تنشأ قبل التعرض للمضاد الحيوي وباستقلال عنه؛ والنسبة المقيسة
نحو $\num{3e-10}$ لكل انقسام خلوي.
\end{solution}
